% !TEX root = statpointmenger.tex
\section{\secC}

Recall that we have for $v,w \in D$
\begin{equation*}
 |v|,|w| \leq |v-w| \leq \tfrac 2 3.
\end{equation*}
Hence,  we obtain for a each curve $\g \in \Cia[0,1](\mathbb R / \mathbb Z)$
with $\E(\g) < \infty $ due to the bi-Lipschitz estimate
\begin{equation*}
 |\g(u+v) - \g(u)| \simeq |v|, \quad \quad
 |\g(u+w) - \g(u)| \simeq |w|, \quad \quad
 |\g(u+v) - \g(u+w)| \simeq |v-w|
\end{equation*}
for all $(u,v,w) \in \mathbb R / \mathbb Z \times D$.
(Here $a\simeq b$ is an abbreviation for the existence of uniform
constants $0<c\le C<\infty$ with $cb\le a\le Cb$.)

We derive the following form for the first variation which is at first site 
much more complicated than the formula derived by Hermes~\cite{hermes}, but due to the special structure of $D$
it is easier to do estimates using this formula.
We abbreviate
\[ R^{p,q}(u_{1},u_{2},u_{3}) := \RR(\g(u_{1}),\g(u_{2}),\g(u_{3})) \]
and we still use
\begin{equation}\tag{\ref{eq:short-notation}}
 \T[\bullet]{v}{w}:=\D[\bullet]vw.
\end{equation}
In contrast to O'Hara's knot energies, we can use a rather direct argument to deduce that
the integral Menger curvature is G\^ateaux differentiable by investigating the integrand, i.e.\@ by looking at the 
Lagrangian
\begin{equation*}
 L(\gamma)(u,v,w) := \frac {\left|\T w0 \wedge \T v0 \right|^q}{|\T w0|^{q}|\T v0|^q |\T vw|^q}|\gamma'(u)| |\gamma'(u+v)| |\gamma'(u+w)|.
\end{equation*}
For $\gamma, h \in \W$ as in the statement of the lemma and $\gamma_\tau := \gamma + \tau h$ one calculates
\begin{align*}
 &\delta L(\gamma;h) (u,v,w) :=\left.  \frac {\partial}{\partial \tau }\left(L_{\gamma_\tau}(u,v,w) \right) \right|_{\tau =0} \\
  &= \Bigg\{ q \left| \T w0 \wedge \T v0 \right|^{q-2}  \frac {\sp{\T w0 \wedge \T v0 ,\T[h] w0 \wedge \T v0 + \T w0 \wedge \T[h] v0 }}{|\T w0|^{p}|\T v0|^p |\T vw|^p} \\
  &- p \frac {\left|\T w0 \wedge \T v0 \right|^q}{|\T w0|^{p+2}|\T v0|^p |\T vw|^p}  \cdot  \sp{\T w0, \T[h] w0} \\
  &- p \frac {\left|\T w0 \wedge \T v0 \right|^q}{|\T w0|^{p}|\T v0|^{p+2} |\T vw|^p}  \cdot  \sp{\T v0, \T[h] v0} \\
  &- p \frac {\left|\T w0 \wedge \T v0 \right|^q}{|\T w0|^{p}|\T v0|^p |\T vw|^{p+2}}  \cdot  \sp{\T vw, \T[h] vw} \\
   &+ R^{p,q}(u, u+w, u+v) \sp{\frac {\dg (u)}{\abs{\dg (u)} }, \frac{\dh(u)}{\abs{\dg (u)} }} \\
  &+ R^{p,q}(u, u+w, u+v) \sp{\frac {\dg (u+v)}{\abs{\dg (u+v)} }, \frac{\dh(u+v)}{\abs{\dg (u+v)} }} \\
  &+ R^{p,q}(u, u+w, u+v) \sp{\frac {\dg (u+w)}{\abs{\dg (u+w)} }, \frac{\dh(u+w)}{\abs{\dg (u+w)} }} \Bigg\} |\dg(u+w)| |\dg(u+v)| |\dg(u)|.
\end{align*}
For future reference, we denote the seven terms one obtains from this formula after factoring out the 
outermost bracket by $\delta L_1, \ldots, \delta L_7$.


\begin{lemma}\label{lem:diff}
 Let \eqref{eq:sub-critical} hold and $\gamma \in \W\rzd$ be an injective regular curve. Then $\E$
 is G\^ateaux differentiable in $\gamma$ and 
 the first variation in direction $h \in \W \rzd$ is given by
 \begin{align}\label{eq:diff}
  \delta \E(\gamma;h) = 6 \iiint_{\mathbb R / \mathbb Z \times D}\delta L(\gamma;h)(u,v,w) \d w \d v \d u
\end{align}
and~\eqref{eq:dE-intro} holds.
% where
%\begin{align*}
% \delta L(\gamma;h) :=&
%  \Bigg(q \left|\T w0 \wedge \T v0\right|^{q-2}  \frac {\sp{\T w0 \wedge \T v0 ,\T[h] w0 \wedge \T v0 + \T w0 \wedge \T[h] v0 }}{|\T w0|^{p}|\T v0|^p |\T vw|^p}   \\ 
%  &- p \frac {\left|\T w0 \wedge \T v0 \right|^q}{|\T w0|^{p-2}|\T v0|^p |\T vw|^p}  \cdot  \sp{\T w0, \T[h] w0} \\
%  &- p \frac {\left|\T w0 \wedge \T v0 \right|^q}{|\T w0|^{p}|\T v0|^{p-2} |\T vw|^p}  \cdot  \sp{\T v0, \T[h] v0} \\
%  &- p \frac {\left|\T w0 \wedge \T v0 \right|^q}{|\T w0|^{p}|\T v0|^p |\T vw|^{p-2}}  \cdot  \sp{\T vw, \T[h] vw} \\
%  &+ R^{p,q}(u, u+w, u+v) \sp{\frac {\dg (u)}{\abs{\dg (u)} }, \frac{\dh(u)}{\abs{\dg (u)} }} \\
%  &+ R^{p,q}(u, u+w, u+v) \sp{\frac {\dg (u+v)}{\abs{\dg (u+v)} }, \frac{\dh(u+v)}{\abs{\dg (u+v)} }} \\
%  &+ R^{p,q}(u, u+w, u+v) \sp{\frac {\dg (u+w)}{\abs{\dg (u+w)} }, \frac{\dh(u+w)}{\abs{\dg (u+w)} }}
% \Bigg)  \\ &|\dg(u+w)| |\dg(u+v)| |\dg(u)| .
% \end{align*}
\end{lemma}

\begin{proof}
Let $U$ be a neighborhood of $\gamma$ in $\W \subset C^{(3p-3)/q-1} \subset C^1$ consisting only of regular curves with
\begin{equation*}
 \inf_{\tilde \gamma \in U, u \in \mathbb R / \mathbb Z} |\tilde \gamma '(u)| =: M_1 >0 
\end{equation*}
and
\begin{equation*}
 \sup_{\tilde \gamma \in U, u\not=v \in \mathbb R / \mathbb Z} \frac {|\tilde \gamma(u) - \tilde \gamma(v)|}{|u-v|} =: M_2 < \infty. 
\end{equation*}
Using 
\begin{multline*}
 \sp{\frac {\T w0} w \wedge \frac{\T v0} v ,\frac {\T[h] w0} w \wedge \frac {\T v0} v 
 + \frac{\T w0} w\wedge \frac {\T[h] v0} v }
 \\= \sp{\left( \frac{\T w0} w - \frac {\T v0}v \right) \wedge \frac{\T v0} v ,
	\left( \frac{\T[h] w0} w - \frac {\T[h] v0}v \right) \wedge \frac{\T v0} v  
	+ \left( \frac{\T w0} w - \frac {\T v0}v \right) \wedge \frac{\T[h] v0} v},
\end{multline*}
\begin{equation*}
 \frac {\T 0w} w \wedge \frac { \T  0v} v = \left( \frac {\T 0w} w- \frac{\T 0v}v \right) \wedge \frac{\T 0v}v,
\end{equation*}
and
\begin{equation*}
 R^{p,q}(u,v,w) = \frac {\left|\left( \frac {\T 0w} w- \frac{\T 0v}v \right) \wedge \T 0v\right|^q}{|w|^{-q} \abs{\T0w}^{p}\abs{\T v0}^{p}\abs{\T vw}^{p}}
\end{equation*}
together with the bi-Lipschitz estimate we infer for $\tilde \gamma \in U$
\begin{equation}  \label{eq:estimateintegrand}
  \left|\delta L(\tilde \gamma; h)(u,v,w) \right|
 \leq C  \frac{|\frac {\T[\tilde \gamma] w0} w - \frac {\T[\tilde \gamma] v0} v|^q \|h'\|_{L^\infty} +
  |\frac {\T[\tilde \gamma] w0} w - \frac {\T[\tilde \gamma] v0} v|^{q-1}|\frac {\T[h] w0} w - \frac {\T[h] v0} v|}{|w|^{p-q} |v|^{p-q} |v-w|^p}.
\end{equation}
So for $0<|\tau|\leq 1$ so small that $\gamma_\tau \in U$
\begin{equation} \label{eq:majorant}
\begin{split}
 &\left|\frac {\partial}{\partial \tau }L(\gamma_\tau)(u,v,w) \right| \\
 &\leq C  \frac{|\frac {\T[\gamma] w0} w - \frac {\T[\gamma] v0} v|^q \|h'\|_{L^\infty} + |\frac {\T[h] w0} w - \frac {\T[h] v0} v|^q \|h'\|_{L^\infty} +
  |\frac {\T[\gamma] w0} w - \frac {\T[\gamma] v0} v|^{q-1}|\frac {\T[h] w0} w - \frac {\T[h] v0} v| + |\frac {\T[h] w0} w - \frac {\T[h] v0} v| ^q }{|w|^{p-q} |v|^{p-q} |v-w|^p}
 \\ &=: g(u,v,w)
 \end{split}
\end{equation}
where $g$ does not depend on $\tau$ and $C$ depends on $M_1, M_2, p, q$ only.

For $f \in \W \rzd$, we have
\begin{equation}\label{eq:esttripleintegral}
\begin{split}
   &\iiint_{\mathbb R / \mathbb Z \times D }\frac{|\frac {\T[f] w0} w - \frac {\T[f] v0} v|^q  }
   {|w|^{p-q} |v|^{p-q} |v-w|^p} \d v \d w \d u \\
  & = \iiint\limits_{\mathbb R / \mathbb Z \times D } 
   \frac{\abs{\int_{0}^{1}\br{f'(u+\theta w)-f'(u+\theta v)}\d\theta}^{q}}{ |w|^{p-q} |v|^{p-q} |w-v|^p} \d w \d v \d u \\
   & \le  C \int_0^{1}\iiint_{\R/\Z\times D } \frac{\abs{f'(u+\theta(w-v))-f'(u)}^{q}}{ |w|^{p-q} |v|^{p-q} |w-v|^p} \d w \d v \d u \d\theta \\
   &\refeq[\le]{intM<seminorm} C\norm f_{{\W}}^{q}.
   \end{split}
\end{equation}
Hence,
\begin{equation}
\begin{split} \label{eq:L1majorant}
&\iiint\limits_{\mathbb R / \mathbb Z \times D }\abs{g(u,v,w)}\d w\d v\d u \\
   &\le C\seminorm \g_{\W[(3p-q-2)/q,q]}^q \|h'\|_{L^\infty} + {}
  \seminorm {h}_{\W[(3p-q-2)/q,q]}^q \br{1+\lnorm[\infty]{h'}}
   \\ & \quad +   C  \left( \iiint\limits_{\mathbb R / \mathbb Z \times D} \frac{|\frac {\T w0} w - \frac {\T v0} v|^q}{|w|^{p-q} |v|^{p-q} |v-w|^p}\d w\d v\d u \right)^{1-\frac 1 q}
   \left( \iiint\limits_{\mathbb R / \mathbb Z \times D} \frac{|\frac {\T[h] w0} w - \frac {\T[h] v0} v|^q}{|w|^{p-q} |v|^{p-q} |v-w|^p} \d w\d v\d u\right)^{\frac 1 q}
  \\
 & \le C \br{\seminorm{\g}_{\W[(3p-2)/q-1,q]}^{q} + \seminorm{h}_{\W[(3p-2)/q-1,q]}^{q}}\br{1+\lnorm[\infty]{h'}}
 +C \seminorm \g_{\W[(3p-q-2)/q,q]}^{q-1}\seminorm {h}_{\W[(3p-q-2)/q,q]}^q.
\end{split}
\end{equation}
So $\delta L(\gamma_\tau; h)$ has a uniform $L^1$-majorant for $\tau$ sufficiently small.
Therefore,
by Lebegue's theorem of dominanted convergence,
we finally can use the fundamental theorem of calculus to write for $\tau$ small
\begin{align*}
 \frac {\E(\gamma+\tau h) - \E(\gamma)}{\tau} 
 &= 6 \iiint\limits_{\mathbb R / \mathbb Z \times D} \int_0^1 \delta L(\gamma_{s \tau};h)(u,v,w) \d s\d u \d v \d w \\
 & \xrightarrow{\tau \rightarrow 0} 6 \iiint\limits_{\mathbb R / \mathbb Z \times D} \delta L (\gamma; h)(u,v,w) \d u \d v \d w.
\end{align*}
Consequently, the first variation exists and has the form~\eqref{eq:diff}.

Using once more the symmetry of the integrand, we can bring this into 
the form~\eqref{eq:dE-intro} as follows.
%
%\begin{corollary}
% Let $\g\in\W\rzd$ be injective and parametrized by arc-length.
% Then, for any $h\in\W\rzd$, the first variation of~$\E$ at~$\g$ in direction~$h$
% can be give by%\vspace{-2ex}
% \begin{equation}\label{eq:dE}
% \begin{split}
%   &\delta\E(\g,h) \\
%   &= \iiint\limits_{(\mathbb R / \mathbb Z)^3} \Bigg\{
% 2q\frac{\abs{\T v0\wedge\T w0}^{q-2}}{\br{\abs{\T vw}\abs{\T v0}\abs{\T w0}}^p}
% \cdot{\sp{\T v0\wedge\T w0,\T v0\wedge\T[h] w0}} \\
%% &\qquad\qquad\quad{}-2q\frac{\abs{\T v0\wedge\T w0}^{q}}{\abs{\T vw}^{p}\abs{\T v0}^{q+2}\abs{\T w0}^{q}}\sp{\T v0,\T[h] v0}\nonumber\\
% &\qquad\qquad\qquad{}-3p\frac{\abs{\T v0\wedge\T w0}^{q}}{\br{\abs{\T vw}\abs{\T v0}\abs{\T w0}}^p}\cdot\frac{\sp{\T vw,\T[h] vw}}{\abs{\T vw}^{2}}\\
% &\qquad\qquad\qquad{}+3\frac{\abs{\T v0\wedge\T w0}^{q}}{\br{\abs{\T vw}\abs{\T v0}\abs{\T w0}}^p}\cdot\sp{\dg,\dh}
%\Bigg\}\d w\d v\d u.%\label{eq:dE}
% \end{split}
% \end{equation}
%\end{corollary}
%
%\begin{proof}
Due to the symmetry of the integrand we have
$L\circ P_{\s}=L$ for any permutation matrix $P_{\s}$, $\in\mathfrak S_{3}$.
%$$
%  L \circ g_{\sigma} = L
%$$
So we obtain
\begin{equation*}
 6\iiint_{\mathbb R/ Z \times D} \delta L(\gamma; h) 
 =  \sum_{\sigma \in\mathfrak S_{3} } \iiint_{P_{\s}(\mathbb R/ Z \times D)} \delta L(\gamma;h)
 = \iiint_{(\mathbb R / \mathbb Z) ^3} \delta L (\gamma;h).
\end{equation*}
The symmetry of $\RR$ now leads to the desired.

Furthermore, by~\eqref{eq:estimateintegrand},
the first variation defines a bounded linear operator on $\W$.
Hence $\E$ is G\^ateaux differentiable.
\end{proof}

In fact, we can even show that $\E$ is $C^1$,
though we will not use this fact in the rest of this article.

\begin{lemma}
 The functional $\E$ is $C^1$ on the subspace of all regular embedded
 $\W$-curves.
\end{lemma}

\begin{proof}
We will prove by contradiction that $\delta \E$
is a continuous map from embedded regular $\W$-curves
to $\br\W^\ast$.
So let us assume that $\delta \E$% : \W \rightarrow (\W)^\ast
was not continuous in $\gamma_{0}$.
Consequently, there are some $\eps_{0}>0$ and sequences $\seqn\g,\seqn h\subset\W$,
$\g_k\to\g_{0}$ in $\W$, $\norm{h_{k}}_{\W}\le1$, with
\begin{equation}\label{eq:contrassumpt}
 |\delta \E(\g_k; h_k) - \E (\gamma;h_k)| \geq \varepsilon_0.
\end{equation}

As in the proof of Lemma~\ref{lem:diff}, we can exploit the embedding $\W \hookrightarrow C^1$ and the openness of the set of regular embedded curves in $C^1$,
to find an open neighborhood $U$ of $\g_{0}$ consisting only of embedded curves, such that
\begin{equation*}
 \inf_{\tilde \gamma \in U, u \in \mathbb R / \mathbb Z} |\tilde \gamma '(u)| =: M_1 >0 
\end{equation*}
and
\begin{equation*}
 \sup_{\tilde \gamma \in U} \|\tilde \gamma\|_{\W} +\sup_{\tilde \gamma \in U, u\not=v \in \mathbb R / \mathbb Z} \frac {|\tilde \gamma(u) - \tilde \gamma(v)|}{|u-v|} =: M_2 < \infty. 
\end{equation*}
After passing to a subsequence
we may assume $\seqn\g\subset U$ and $h_{k}\to h_{0}\in\W$ in $C^{1}$
due to the compactness of the embedding $\W \hookrightarrow C^1$
which then also gives
\begin{equation*}
 \delta L(\g_k;h_k) - \delta L(\g_0;h_0) \rightarrow 0
\end{equation*}
pointwise almost everywhere on $\R/\Z\times D$ and hence in measure, i.e., for all $\varepsilon >0$ we have
\begin{equation}\label{eq:lebto0}
 \lim_{k \rightarrow \infty} \mathcal L^3 \br{A_{\eps,k}}=0,
\end{equation}
where $\mathcal L^3$ denotes the Lebesgue measure
and
\begin{equation*}\delimitershortfall-1pt
 A_{\eps,k}:= \sett{(u,v,w) \in \mathbb R / \mathbb Z \times D}
 {\abs{\delta L(\g_k;h_k)(u,v,w)-\delta L(\g_0;h_0)(u,v,w)} \geq \varepsilon}. 
\end{equation*}

For all $\varepsilon >0$ we can deduce from \eqref{eq:estimateintegrand}
using Young's inequality that there is a $C_{\eps}>0$ with
\begin{equation}
\begin{split}
 \left|\delta L(\gamma_{k};h_{k})(u,v,w) \right|
 &\leq   C_\varepsilon  \frac{|\frac {\T[\g_k] w0} w - \frac {\T[\g_k] v0} v|^q }{|w|^{p-q} |v|^{p-q} |v-w|^p} 
  + \varepsilon \frac { |\frac {\T[h_{k}] w0} w - \frac {\T[h_{k}] v0} v|^q}{|w|^{p-q} |v|^{p-q} |v-w|^p} 
  \\ &=: C_{\eps}g^{(1)}_k(u,v,w) + \varepsilon g^{(2)}_k(u,v,w)
 \end{split} 
\end{equation}
for all $k\in\N\cup\set0$.
Since the first summand converges in $L^1$ as $k\to\infty$, it is uniformly integrable,
so % i.e., for $\eta :=\frac\eps{C_{\eps}}>0$
there is a $\delta_\eps >0$ such that
$\mathcal L^{3} (E) \leq \delta_\eps$ for any measurable subset $E \subset \mathbb R / \mathbb Z \times D$ implies
\begin{equation}
 \iiint_E g_k^{(1)}  \leq \frac\eps{C_{\eps}} \qquad\text{for all }k\in\N\cup\set0.
\end{equation}
Furthermore we infer from \eqref{eq:intM<seminorm}
\begin{equation*}
 \iiint_{\mathbb R / \mathbb Z \times D} g^{(2)}_k \d u\d v\d w \leq C \norm {h_{k}}_{\W} \leq C.
\end{equation*}
By~\eqref{eq:lebto0} there is some $k_0=k_{0}(\eps) \in \mathbb N$ with
\begin{equation*}
 \mathcal L^3 \br{A_{\eps,k}} \leq \delta_\eps
 \qquad\text{for all } k \geq k_0
\end{equation*}
which yields for
$k \geq k_0$ and
$B_{\eps,k} := \mathbb R / \mathbb Z \times D \setminus A_{\eps,k}$
\begin{align*}
 &\tfrac16|\delta \E(\g_k; h_k) - \delta\E (\gamma_{0};h_{0})| \\
 &\leq \iiint_{A_{\eps,k}} |\delta L(\g_k;h_k)-\delta L(\g_0;h_0) |
  + \iiint_{B_{\eps,k}} |\delta L(\g_k;h_k)-\delta L(\g_0;h_0) | \\
 & \leq \iiint_{A_{\eps,k}} \left(C_{\eps} g^{(1)}_k +  \varepsilon g^{(2)}_k \right)
 + \iiint_{A_{\eps,k}} \left(C_{\eps} g_{0}^{(1)} +  \varepsilon g_{0}^{(2)} \right)
 + \mathcal L^3(B_{\eps,k}) \varepsilon \\
 & %\leq 2C_{\eps}\eta+2C\eps+\mathcal L^3(B) \varepsilon
  \leq C\eps.
\end{align*}
Hence,
\begin{align*}
 \eps_{0}&\refeq[\le]{contrassumpt}\abs{\delta \E(\g_k; h_k) - \E (\gamma;h_k)} \\
 &\le \abs{\delta \E(\g_k; h_k) - \E (\gamma;h)} + \abs{\delta \E(\g; h_k) - \E (\gamma;h)} \\
 &\le C\eps + C\norm{h_{k}-h}_{\W}
\end{align*}
for all $\varepsilon >0$ and $k\ge k_{0}(\eps)$ which leads to a contradiction.
\end{proof}


\section{\secD}

For the rest of this section, let us restrict to the case that $ \g$ is parametrized by arc-length.
Then we get using Lemma \ref{lem:diff} 
\begin{align*}
 \delta \E (\gamma;h) := 6q \tilde Q^{p,q} (\gamma,h) + 6 R^{p,q}_1 (\gamma h)
\end{align*}
where
\begin{multline*}
 \tilde Q^{p,q}(\gamma,h):= \iiint_{\mathbb R / \mathbb Z \times D} \left|\T w0 \wedge \T v0\right|^{q-2}  \frac {\sp{\T w0 \wedge \T v0 ,\T[h] w0 \wedge \T v0 + \T w0 \wedge \T[h] v0 }}{|\T w0|^{p}|\T v0|^p |\T vw|^p}   
\d w \d v \d u
\end{multline*}
and \begin{align*}
  R_1 (\gamma,h) := \iiint_{\mathbb R / \mathbb Z \times D} \Bigg(
  &- p \frac {\left|\T w0 \wedge \T v0 \right|^q}{|\T w0|^{p+2}|\T v0|^p |\T vw|^p}  \cdot  \sp{\T w0, \T[h] w0} \\
  &- p \frac {\left|\T w0 \wedge \T v0 \right|^q}{|\T w0|^{p}|\T v0|^{p+2} |\T vw|^p}  \cdot  \sp{\T v0, \T[h] v0} \\
  &- p \frac {\left|\T w0 \wedge \T v0 \right|^q}{|\T w0|^{p}|\T v0|^p |\T vw|^{p+2}}  \cdot  \sp{\T vw, \T[h] vw} \\
  &+ R^{p,q}(u, u+w, u+v) \sp{\dg (u), \dh(u)} \\
  &+ R^{p,q}(u, u+w, u+v) \sp{\dg (u+v),\dh(u+v)} \\
  &+ R^{p,q}(u, u+w, u+v) \sp{\dg (u+w),\dh(u+w)}
 \Bigg)\d w \d v \d u.
\end{align*}

For $q=2$ will see that $\tilde Q^p := \tilde Q^{p,2}$ contains the highest order term of the Euler-Lagrange operator.
To see this we use
$$
 \left\langle a\wedge b , c\wedge d  \right\rangle = \det \begin{pmatrix}\sp{a,c} & \sp{a,d} \\
                                                           \sp{b,c} & \sp{b,d}
                                                          \end{pmatrix} 
$$
to get 
\begin{equation}\label{eq:abac}
 \left\langle a \wedge b, a \wedge c\right\rangle = \sp{a,a}\sp{c,b} - \sp{a,c}\sp{a,b} 
 = |a|^2\left\langle P^\bot_{a} b, c \right\rangle.
\end{equation}
Hence,
\begin{align*}
\tilde Q^{p}(\gamma;h) &= \iiint\limits_{\mathbb R / \mathbb Z \times D } \Bigg(
 \frac{\sp{P^\bot_{\T v0}\T w0,\T[h] w0}}{\abs{\T vw}^p\abs{\T v0}^{p-2}\abs{\T w0}^p}
+
 \frac{\sp{P^\bot_{\T w0}\T v0,\T[h] v0}}{\abs{\T wv}^p\abs{\T v0}^p\abs{\T w0}^{p-2}}
 \Bigg)\d w \d v \d u
\\
&=    \iiint\limits_{\mathbb R / \mathbb Z \times D } \Bigg( \frac{\sp{P^\bot_{\T v0}\T w0,\T[h] w0}}{\abs{v-w}^p\abs{v}^{p-2}\abs{w}^p}
+
 \frac{\sp{P^\bot_{\T w0}\T v0,\T[h] v0}}{\abs{v-w}^p\abs{v}^p\abs{w}^{p-2}}
 \Bigg)\d w \d v \d u + R_2(\gamma;h) \\
& =   \iiint\limits_{\mathbb R / \mathbb Z \times D } \Bigg( \frac{\sp{P^\bot_{\T v0} \br{\frac {\T w0}{w}
- \frac {\T v0}{v}},\frac{\T[h] w0}{w}}}{\abs{v-w}^p\abs{v}^{p-2}\abs{w}^{p-2}}
-
 \frac{\sp{P^\bot_{\T w0} \br{\frac {\T w0}{w}
- \frac {\T v0}{v}},\frac{\T[h] v0}{v}}}{\abs{v-w}^p\abs{v}^{p-2}\abs{w}^{p-2}}
 \Bigg)\d w \d v \d u + R_2(\gamma;h) \\
%  &=   \iiint\limits_{\mathbb R / \mathbb Z \times D } \Bigg( \frac{\sp{P^\bot_{\T v0}\T w0,\T[h] w0}}{\abs{v-w}^p\abs{v}^{p-2}\abs{w}^p}
% +
%  \frac{\sp{P^\bot_{\T w0}\T v0,\T[h] v0}}{\abs{v-w}^p\abs{v}^p\abs{w}^{p-2}}
%  \Bigg)\d w \d v \d u + R_2(\gamma;h) \\
& =   \iiint\limits_{\mathbb R / \mathbb Z \times D } \Bigg( \frac{\sp{\frac {\T w0}{w}
- \frac {\T v0}{v},\frac{\T[h] w0}{w}-\frac{\T[h] v0}{v}}}{\abs{v-w}^p\abs{v}^{p-2}\abs{w}^{p-2}}
 \Bigg)\d w \d v \d u + R_2(\gamma;h) - R_3(\gamma;h) \\
\end{align*}
where
\begin{align*}
 R_2 (\gamma;h) := \iiint\limits_{\mathbb R / \mathbb Z \times D } \Bigg(
 &\sp{P^\bot_{\T v0}\T w0,\T[h] w0}  \left( \frac 1{\abs{\T vw}^p\abs{\T v0}^{p-2}\abs{\T w0}^p} 
 - \frac 1{\abs{v-w}^p\abs{v}^{p-2}\abs{w}^p}\right)
 \\ &+
 \sp{P^\bot_{\T w0}\T v0,\T[h] v0} \left( \frac 1 {\abs{\T wv}^p\abs{\T v0}^p\abs{\T w0}^{p-2}}
  - \frac 1 {\abs{v-w}^p\abs{v}^p\abs{w}^{p-2}}\right)
 \Bigg)\d w \d v \d u
\end{align*}
and
\begin{align*}
 R_3 (\gamma;h) :=   \iiint\limits_{\mathbb R / \mathbb Z \times D } \Bigg( \frac{\sp{P^{\top}_{\T v0}
\left( \frac{ \T w0 } w - \frac{ \T v0 } v \right),\frac{\T[h] w0}w}}{\abs{v-w}^p\abs{v}^{p-2}\abs{w}^{p-2}}
+
 \frac{\sp{P^{\top}_{\T w0}\left(\frac {\T v0}{v} -\frac {\T w0}{w}\right) ,\frac{\T[h] v0}v}}{\abs{v-w}^p\abs{v}^{p-2}\abs{w}^{p-2}}
 \Bigg)\d w \d v \d u.
\end{align*}
Using
\begin{align*}
 Q^{(p)}(\gamma;h) :=   \iiint\limits_{\mathbb R / \mathbb Z \times D } \Bigg( \frac{\sp{\frac {\T w0}{w}
- \frac {\T v0}{v},\frac{\T[h] w0}{w}-\frac{\T[h] v0}{v}}}{\abs{v-w}^p\abs{v}^{p-2}\abs{w}^{p-2}}
 \Bigg)\d w \d v \d u
\end{align*}
we hence get
\begin{equation}\label{eq:Re}
 \delta M^{p,2} = 12 \left( Q^{(p)} + \tfrac12R_1 +R_2 + R_3  \right).
\end{equation}


\begin{proposition}\label{prop:Q}
 The functional $\Q$ is bilinear on $\br{\W[3p/2-2,2]}^{2}$, more precisely
 \[ \Q(f,g) = \sum_{k\in\Z} \rho_k\sp{\hat f_k,\hat g_k}_{\C^d} \qquad \text{where } \rho_k = c\abs k^{3p-4} + o\br{\abs k^{3p-4}} \text{ as }\abs k\nearrow\infty \]
 and $c>0$.
 Here $\hat\cdot_k$ denotes the $k$-th Fourier coefficient
 \[
  \hat f_k := \int_{0}^1 f(x) e^{-2\pi ikx } \d x.
 \]
\end{proposition}

\begin{proof}
   Testing with the basis ${\rm e}_l \cdot e^{2\pi i k x}$ of $L^2$, $l=1,2,3$, $k\in \mathbb Z$, where ${\rm e}_1, {\rm e}_2, {\rm e}_3$ is the standard basis of $\mathbb R ^3$,
   we get
\begin{equation*}
      \Q(f,g) =
       \sum_{k\in\Z} \sp{\hat f_k,\hat g_k}_{\C^d} \rho_k
 \end{equation*}
 where
 \begin{equation*} 
       \rho_k :=  
       \iint\limits_{D} 
       \frac{ \abs{ \frac {e^{2\pi i  kw} -1} {w} - \frac {e^{2\pi i k v} -1} {v} }^2}{|v|^{p-2}|w|^{p-2}|v-w|^{p}} \d w\d v.
  \end{equation*}
 A simple substitution leads to 
 \begin{align*}
  \rho_k = \abs k^{3p - 4}
  \iint\limits_{D_k} 
       \frac{ \abs{ \frac {e^{2\pi i  w} -1} {w} - \frac {e^{2\pi i  v} -1} {v}}^2}{|v|^{p-2}|w|^{p-2}|v-w|^{p}} \d w\d v 
 \end{align*}
 where $D_k := k \cdot D.$ We use the fundamental thorem of calculus and Jensen's inequality to get
 \begin{align*}
  \frac{ \abs{ \frac {e^{2\pi i  w} -1} {w} - \frac {e^{2\pi i  v} -1} {v}}^2}{|v|^{p-2}|w|^{p-2}|v-w|^{p}}
  &\leq 4\pi^{2}\int_0^1 \frac {|e^{2\pi i \theta w} - e^{2\pi i \theta v} |^2} {|w|^{p-2} |v|^{p-2} |v-w|^{p}} \d \theta\\
  &= 4\pi^{2} \int_0^1 \frac {|e^{2\pi i \theta (v-w)} - 1 |^2} {|w|^{p-2} |v|^{p-2} |v-w|^{p}} \d \theta.
 \end{align*}
 Hence, 
 \begin{align*}
  \iint\limits_{v<0, w>0} \frac{ \abs{ \frac {e^{2\pi i k w} -1} {w} - \frac {e^{2\pi i k v} -1} {v}}^2}{|v|^{p-2}|w|^{p-2}|v-w|^{p}}
  &\leq C\iint\limits_{v<0, w>0} \int_0^1 \frac {|e^{2\pi i \theta (v-w)} - 1 |^2} {|w|^{p-2} |v|^{p-2} |v-w|^{p}} \d \theta \d w \d v
  \\ &= C\left( \int_0^1  {\theta^{3p-6}} \d \theta \right) \iint\limits_{v<0, w>0}  \frac {|e^{2\pi i  (v-w)} - 1 |^2} {|w|^{p-2} |v|^{p-2} |v-w|^{p}}  \d w \d v
  \\ &\leq C \iint\limits_{v<0, w>0}  \frac {|e^{2\pi i  (v-w)} - 1 |^2} {|w|^{p-2} |v|^{p-2} |v-w|^{p}}  \d w \d v
  \\ & = C \int_0^\infty \int_{0}^{\tilde w} \frac {|e^{-2\pi i \tilde w} - 1 |^2} {|\tilde w -v|^{p-2} |v|^{p-2} |\tilde w|^{p}}  \d v \d \tilde w
  \\ & = C  \int_0^\infty \int_0^1  \frac {|e^{-2\pi i \tilde w} - 1 |^2} {|1 -t|^{p-2} |t|^{p-2} |\tilde w|^{3p-5}}  \d t \d \tilde w
  \\ & \le C \int_0^\infty \frac {1-\cos2\pi\tilde w} {|\tilde w|^{3p-5}} \d \tilde w < \infty.
 \end{align*}
 Thus, we have shown that
 \begin{equation*}
  \frac {\rho_k} {|k|^{3k-4}} \xrightarrow{|k| \rightarrow \infty} \iint_{v<0, w>0} \frac{ \abs{ \frac {e^{2\pi i k w} -1} {w} - \frac {e^{2\pi i k v} -1} {v}}^2}{|v|^{p-2}|w|^{p-2}|v-w|^{p}}\d w\d v \in (0, \infty).
 \end{equation*}
\end{proof}

In the following statement, we use the symbol $\otimes$ to
denote any kind of product structure, such as cross product, dot product,
scalar or matrix multiplication.

\begin{lemma}
 The term $R^{(p)} := \tfrac12R_1 + R_2 + R_3$ is a finite sum of terms of the form
 $$
  \iiint\limits_{\mathbb R / \mathbb Z \times D} \idotsint\limits_{[0,1]^K} g^p(u,v,w;s_{1},\dots,s_{K-2}) \otimes h'(u+s_{K-1}v + s_K w) 
  \d\theta_{1}\cdots\d\theta_{K}\d v \d w \d u
 $$
 where $\G : (0,\infty)^3 \rightarrow \mathbb R$ is an analytic funtion,
 $s_{j}\in\set{0,\theta_{j}}$ for $j=1,\dots,K$,
 \begin{multline*}
  g^p(u,v,w;s_{1},\dots,s_{K-2})=%\iiint\limits_{\mathbb R / \mathbb Z \times D} \int_{[0,1]^K}
   \G \left(\frac{|\T 0w|}{|w|}, \frac{|\T 0v|}{|v|}, \frac{|\T vw|}{|v-w|}\right) 
  \Gamma(u,v,w,s_{1},s_{2})\cdot{}\\
  {}\cdot\left( \bigotimes_{i=3}^{K_1} \dg(u+ s_i v) \right) \otimes \left(\bigotimes_{j=K_1}^{K_2} \dg(u + s_i w) \right) \otimes \left(\bigotimes_{j=K_2}^{K-2} \dg(u+v + s_i (w-v)) \right),
%  \\ \d s_1 \cdots \d s_K \d w \d v,
 \end{multline*}
 and
 $\Gamma(u,v,w,s_{1},s_{2})$ is a term of one of the four types
 \begin{align*}
  &\frac{\br{\gamma'(u+s_1 w) - \gamma'(u+s_1 v)}\otimes\br{\gamma'(u+s_2 w) - \gamma'(u+s_2 v)}}{|v|^{p-2} |w|^{p-2} |v-w|^{p}}, \\
  &\frac{|\gamma'(u+s_1 w) - \gamma'(u+s_2 w)|^2}{|v|^{p-2} |w|^{p-2} |v-w|^{p}}, \\ 
  & \frac{|\gamma'(u+s_1 v) - \gamma'(u+s_2 v)|^2}{|v|^{p-2} |w|^{p-2} |v-w|^{p}}, \\ 
  & \frac{|\gamma'(u+v + s_1 (w-v)) - \gamma'(u+v+s_2 (w-v))|^2}{|v|^{p-2} |w|^{p-2} |v-w|^{p}}.
 \end{align*}
\end{lemma}

\begin{proof}
 Using
 \begin{align*}
  \frac{\left|\T w0 \wedge \T v0 \right|^2}{|\T w0|^{p+2}|\T v0|^p |\T vw|^p}  \cdot  \sp{\T w0, \T[h] w0}
  = \G \left(\frac{|\T 0w|}{|w|}, \frac{|\T 0v|}{|v|}, \frac{|\T vw|}{|v-w|}\right)
 \frac{\left|\T w0 \wedge \T v0 \right|^2}{|w|^{p+2}|v|^p |v-w|^p}  \cdot  \sp{\T w0, \T[h] w0} 
 \end{align*}
 where
$
  \G (z_1, z_2, z_2) = z_1^{-p-2} z_2 ^{-p} z_3 ^{-p}
$
together with
 \begin{equation*}
  \frac {\T w0} w \wedge \frac {\T v0} v
  = \left(\frac {\T w0} w- \frac {\T v0} v \right)\wedge \frac {\T v0} v
  = \int_{0}^{1} \int_{0}^1 \left(\gamma'(u + s_1 w) - \gamma'(u+s_1 v) \right) \wedge \gamma'(u+s_3 v) \d s_1 \d s_3
 \end{equation*}
 we see that the first term of $R^1$ is of type 1. Similarly, one gets that all the terms of $R^1$ are
of type 1.


For the term $R^2$ we use
\begin{multline*}
 \sp{P^\bot_{\T 0v} (\T 0w), \T[h] 0w}
 =\frac 1 {|\T 0v|^2} \sp{(\T 0v \wedge \T 0w), (\T 0v \wedge \T[h] 0w)}  \\
 \refeq{abac} \frac {v^{2}w^{2}} {|\T 0v|^2} \idotsint_{[0,1]^4} \gamma'(u+s_1 v) \otimes \gamma'(u+s_2 w) \otimes \gamma'(u+s_3 v) \otimes 
 h'(u+s_4 w) \d s_1 \d s_2 \d s_3 \d s_4
\end{multline*}
together with the fact that for  $w \in \mathbb R$, $u \in \mathbb R / \mathbb  Z$
 \begin{equation} \label{eq:OHaraTypeTerm}
 \begin{split}
  \frac 1{|\T w0|^\alpha} - \frac 1 {|w|^\alpha} &= 2\G[\alpha]\br{\frac{\abs{\T w0}}{|w|}} \frac {1-\frac {\abs{\T w0}^2}{|w|^2}} {|w|^\alpha}\\
  &= \G[\alpha] \br{\frac{\abs{\T w0}}{|w|}} \iint\limits_{[0,1]^2} \frac {|\gamma'(u+s_1 w) - \gamma'(u+s_2 w)|^2}{|w|^\alpha} \d s_1 \d s_2
  \end{split}
  \end{equation}
 where $\G[\alpha] (z) = \tfrac12\cdot\frac {1-z^\alpha}{1-z^2} \cdot z^{-\alpha}$ is analytic on $(0,\infty)$ for $\alpha>0$.
 Both equations together with
\begin{align*}
 \frac 1{\abs{\T vw}^p\abs{\T v0}^{p-2}\abs{\T w0}^p} 
 &- \frac 1{\abs{v-w}^p\abs{v}^{p-2}\abs{w}^p}
\\ &= \left(\frac 1{\abs{\T vw}^p\abs{\T v0}^{p-2}\abs{\T w0}^p} 
 - \frac 1{\abs{\T vw}^p\abs{\T v0}^{p-2}\abs{w}^p} \right)
\\ &+ \left(\frac 1{\abs{\T vw}^p\abs{\T v0}^{p-2}\abs{w}^p} 
 - \frac 1{\abs{\T vw}^p\abs{v}^{p-2}\abs{w}^p} \right)
\\ &+ \left(\frac 1{\abs{\T vw}^p\abs{v}^{p-2}\abs{w}^p} 
 - \frac 1{\abs{v-w}^p\abs{v}^{p-2}\abs{w}^p} \right)
\end{align*}
show that the frist term of $R_2$ is the sum of terms of type 2 to 4. Similarly for the second term
in $R_2$.

Let us turn to the last term, $R_3$.
Again, we restrict to the first term, the second is parallel.
We obtain
\begin{align*}
 \frac{\sp{P^{\top}_{\T v0}
\left( \frac{ \T w0 } w - \frac{ \T v0 } v \right),\frac{\T[h] w0}w}}{\abs{v-w}^p\abs{v}^{p-2}\abs{w}^{p-2}}
&= \frac{\sp{\frac{ \T w0 } w - \frac{ \T v0 } v,\fracabs{\T v0}}}{\abs{v-w}^p\abs{v}^{p-2}\abs{w}^{p-2}}\sp{\fracabs{\T v0} ,\frac{\T[h] w0}w} \\
&= \abs{\frac{\T v0}v}^{-2}\frac{\sp{\frac{ \T w0 } w - \frac{ \T v0 } v,\frac{\T v0}v}}{\abs{v-w}^p\abs{v}^{p-2}\abs{w}^{p-2}} \sp{\frac{\T v0}v ,\frac{\T[h] w0}w}\\
&= \abs{\frac{\T v0}v}^{-2}\frac{\br{\sp{\frac{ \T w0 } w , \frac{ \T v0 } v}-1}}{\abs{v-w}^p\abs{v}^{p-2}\abs{w}^{p-2}} \sp{\frac{\T v0}v ,\frac{\T[h] w0}w} \\
&= -\tfrac12\abs{\frac{\T v0}v}^{-2}\frac{{\abs{\frac{ \T w0 } w - \frac{ \T v0 } v}^{2}}}{\abs{v-w}^p\abs{v}^{p-2}\abs{w}^{p-2}} \sp{\frac{\T v0}v ,\frac{\T[h] w0}w},
\end{align*}
which gives rise to type~4.
\end{proof}


 Our next task is to show that
  $R^p$
 is in fact a lower-order term. More precisely, we have
 
 \begin{proposition}[Regularity of the remainder term]\label{prop:reg-rem}
  If $\g\in\Wia[(3p-4)/2+\s,2]\rzd$ for some $\s\ge0$ then
  $R^p(\g,\cdot)\in\br{\W[3/2-\s+\eps,2]}^{*}$ for any $\eps>0$.
 \end{proposition}
 
 This statement together with Proposition~\ref{prop:Q} immediately leads to
 the proof of the regularity theorem which is deferred to the end of
 this section.
 
 To prove Proposition~\ref{prop:reg-rem}% for such a term
 , we first note that,
 by partial integration, the terms of $\Re[](\g,h)$ may be transformed into
\begin{align*}
  &\idotsint_{[0,1]^{K}}\iint_{D}\int_{\R/\Z}
   \br{(-\Delta)^{\tilde\s/2}\gp(\cdot,v,w;s_{1},\dots,s_{K-2})}(u) \otimes\br{(-\Delta)^{-\tilde\s/2}\dh}(u+s_{K-1}v+s_{K}w)
   \d u\d v\d w \d\theta_{1}\cdots\d\theta_{K} \\
  &\le \idotsint_{[0,1]^{K}}
  \iint_{D}\norm{\gp (\cdot,v, w;\dots)}_{\W[\tilde\s,1]\rzd}\d v\d w
  \d\theta_{1}\cdots\d\theta_{K-2} \norm{(-\Delta)^{-\tilde\s/2}\dh}_{L^{\infty}\rzd} \\
 &\le \idotsint_{[0,1]^{K}}
  \iint_{D}\norm{\gp (\dots)}_{\W[\tilde\s,1]\rzd}\d v \d w
  \d\theta_{1}\cdots\d\theta_{K-2}\norm{h}_{\W[3/2+\eps/2-\tilde\s,2]\rzd},
 \end{align*}
  where $\tilde\s\in\R$, $\eps>0$ can be chosen arbitrarily, and $(-\Delta)^{\tilde\s/2}$ denotes the fractional
  Laplacian.
  We let $\tilde\s:=0$ if $\s=0$ and $\tilde\s:=\s-\frac\eps2$ otherwise.
  Now the claim directly follows from the succeeding auxiliary result.
 
 \begin{lemma}[Regularity of the remainder integrand]\label{lem:reg-int}
  Let $\g\in\Wia[(3p-4)/2+\s,2]$.
  \begin{itemize}
  \item If $\s=0$ then $\gp\in L^{1}(\R/\Z\times D,\R^{n})$ and
  \item if $\s>0$ then $((v,w) \mapsto \gp(\cdot,v,w;\dots))\in L^{1}(D,\W[\tilde\s,1](\R/\Z, \R^{n}))$ for any $\tilde\s<\s$.
  \end{itemize}
  The respective norms are bounded independently of $s_{1},\dots,s_{K}$.
 \end{lemma}
 
 \begin{proof}
  Recall that the argument of $\G$ is compact and bounded away from zero.
  Using arc-length parametrization, we immediately obtain for the first type
  \begin{align*}
   \iiint\limits_{\mathbb R / \mathbb Z \times D }\abs{\gp(u,v,w)} \d v \d w \d u &\le C \iiint\limits_{\mathbb R / \mathbb Z \times D } 
   \frac{\abs{\dg(u+s_{1}w)-\dg(u+s_{1}v)}^{2}}{ |w|^{p-2} |v|^{p-2} |w-v|^p} \d w \d v \d u \\
   & =  C \iiint\limits_{\R/\Z \times D } \frac{\abs{\dg(u+s_{1}(w-v))-\dg(u)}^{2}}{ |w|^{p-2} |v|^{p-2} |w-v|^p} \d w \d v \d u \\
   &\refeq[\le]{intM<seminorm} C\norm\g_{\W[(3p-4)/2,2]}^{2}.
   \end{align*}
  For a term of the second type we get
  \begin{align*}
   \iiint\limits_{\mathbb R / \mathbb Z \times D } & \abs{\gp(u,v,w)} \d v \d w \d u \le C \iiint\limits_{\mathbb R / \mathbb Z \times D } 
   \frac{\abs{\dg(u+s_{1}w)-\dg(u+s_{2}w)}^{2}}{ |w|^{p-2} |v|^{p-2} |w-v|^p} \d w \d v \d u \\
   & =  C \int_{\R/\Z} \int_0^{2/3} \frac{\abs{\dg(u+s_{1}w)-\dg(u+s_2 w)}^{2}}{ |w|^{p-2}} \left(\int_0^{2/3-w} \frac 1 {v^{p-2}(v+w)^{p}}\d v\right) \d w \d u \\
   & \le C \int_{\R/\Z} \int_0^{2/3} \frac{\abs{\dg(u+s_{1}w)-\dg(u+s_2 w)}^{2}}{ |w|^{3p-5}} \left(\int_{0}^{\infty} \frac 1 {t^{p-2}(1+t)^{p}}\d t\right) \d w \d u \\
   & \le  C\seminorm\g_{\W[(3p-4)/2,2]}
  \end{align*}
  and of course the same estimate is true for a term of the third kind. For term of type four, we get along the same lines
  \begin{align*}
   \iiint\limits_{\mathbb R / \mathbb Z \times D } & \abs{\gp(u,v,w)} \d v \d w \d u \le C \iiint\limits_{\mathbb R / \mathbb Z \times D } 
   \frac{\abs{\dg(u+v+s_{1}(w-v)-\dg(u+v + s_{2}(w-v))}^{2}}{ |w|^{p-2} |v|^{p-2} |w-v|^p} \d w \d v \d u \\
   & =  C \int_{\rzd} \int_0^{2/3} \frac{\abs{\dg(u+s_{1}(w+v))-\dg(u+s_2 (w+v))}^{2}}{ |w|^{p-2}} \left(\int_{0}^{2/3- w} \frac 1 {(v+w)^{p} v^{p-2}}\d v\right) \d w \d u \\
   & \le C \int_{\rzd} \int_0^{2/3} \frac{\abs{\dg(u+s_{1}\tilde w)-\dg(u+s_2 \tilde w)}^{2}}{ |\tilde w|^{3p-5}} \left(\int_{0}^{\infty} \frac 1 {(t +1)^{p}t^{p-2}}\d t\right) \d w \d u \\
   & \le  C\seminorm\g_{\W[(3p-4)/2,2]}.
  \end{align*}
  
  We prove the second claim only for terms of the first type, as the arguments for all other terms follow the same line of argments.
  We will derive a suitable bound on $\norm{\gp(\cdot,w)}_{\W[\tilde\s,r]}$ for some $r>1$.
  To this end, we choose $q_{1},\dots,q_{K-2}$, which will be determined more precisely later on,
  such that
  \begin{equation*}
    \sum_{i=1}^{K-2} \frac 1 {q_i} = \frac 1 r.
  \end{equation*}
  The product rule, Lemma~\ref{lem:product}, then leads to 
  \begin{align*} 
  \norm{\gp(\cdot,w)}_{\W[\tilde\s,r]} &\le C\norm{\G}_{\W[\tilde\s,q_{1}]}\frac{\norm{\dg(\cdot+s_{1}w)-\dg(\cdot+s_{1}v)}_{\W[\tilde\s,2q_{2}]}^{2}}{\abs v ^{p-2}\abs w^{p-2} \abs{w-v}^p}{\prod_{i=3}^{K-2}\norm\dg_{\W[\tilde\s,q_{i}]}}. 
  \\ & =C\norm{\G}_{\W[\tilde\s,q_{1}]}\frac{\norm{\dg(\cdot+s_{1}(w-v))-\dg(\cdot)}_{\W[\tilde\s,2q_{2}]}^{2}}{\abs v ^{p-2}\abs w^{p-2} \abs{w-v}^p}{\prod_{i=3}^{K-2}\norm\dg_{\W[\tilde\s,q_{i}]}}. 
  \end{align*}
  
  For the second factor, we now choose $q_{2}>r$ so small that
  $\W[\s,2]$ embeds into $\W[\tilde\s,2q_{2}]$.
  To this end, we set $\frac1r:=1-(\s-\tilde\s)$ and $\frac1{q_{2}}:=1-2(\s-\tilde\s)$.
  and $q_{i}:=\frac{K-3}{\s-\tilde\s}$ for $i=1,3,4, \ldots, K-2$.
  
  Then for the first factor we apply the chain rule, Lemma~\ref{lem:chain}.
  Recall that $\G$ is analytic and its argument is bounded below away from zero
  and above by~$1$.
  We infer
  \begin{equation*}
   \norm{\G}_{\W[\tilde\s,q_{1}]}
   \leq C \norm\dg_{\W[ \tilde \sigma ,q_{1}]}.
  \end{equation*}
  The Sobolev embedding gives
  \[ \norm{\dg}_{\W[ \tilde \sigma ,q_{i}]} \le C \norm\g_{\W[(3p-4)/2+\s,2]} \le  C \qquad \text{for } i=1,3,4,\dots,K-2. \]
  Together this leads to
  \[ \norm{\gp(\cdot,v,w)}_{\W[\tilde\s,r]} \le C\frac{\norm{\dg(\cdot+s_{1}(w-v))-\dg(\cdot)}_{\W[\s,2]}^{2}}{\abs v^{p-2}\abs w^{p-2} \abs{w-v}^p} \]
  and finally
  \begin{align*}
     \iint_{D} \norm{\gp(\cdot,w)}_{\W[\tilde\s,r]} \d v \d w
     &\le C \iint_{D} \frac{\norm{\dg(\cdot+s_{1}(w-v))-\dg(\cdot)}_{\W[\s,2]}^{2}}{\abs v^{p-2}\abs w^{p-2} \abs{w-v}^p} \d v \d w \\
     &\refeq[\le]{intM<seminorm} C\norm\g_{{\W[(3p-4)/2,2]}}.
  \end{align*} 
 \end{proof}
 
 \begin{proof}[Theorem~\ref{thm:smooth}]
  We start with the Euler-Lagrange Equation
  \begin{equation}\label{eq:elg}
   \delta\E(\g,h)+\lambda\sp{\dg,\dh}_{L^2} = 0
  \end{equation}
  for any $h\in C^\infty\rzd$
  where $\lambda\in\R$ is a Lagrange parameter stemming from the side condition (fixed length).
  Using~\eqref{eq:Re} this reads
  \begin{equation}\label{eq:elg2}
   12\Q(\g,h) + \lambda\sp{\dg,\dh}_{L^2} + 12\Re[](\g,h) = 0.
  \end{equation}
  Since first variation of the length functional satisfies
  \begin{equation*}
   \sp{\dg,\dh}_{L^2} = \sum_{k\in \mathbb Z} |2\pi k|^2 \sp{\hat\g_{k},\hat h_{k}}_{\C^{d}},
  \end{equation*}
  we get using  Proposition~\ref{prop:Q} that there is a $\tilde c>0$ such that
  \begin{equation}\label{eq:LinearTerm}
  12\Q(\g,h) + \lambda\sp{\dg,\dh}_{L^2} = \sum_{k\in\Z} \tilde\rho_k\sp{\hat \g_k,\hat h_k}_{\C^d}
  \end{equation}
  where
  \[ \tilde\rho_k = \tilde c\abs k^{3p-4} + o\br{\abs k^{p-1}}
     \qquad\text{as }\abs k\nearrow\infty. \]
  Assuming that $\g\in\Wia[(3p-4)/2+\s,2]$ for some $\s\ge0$,
  we infer
  \[ 12\Q(\g,\cdot) + \lambda\sp{\dg,\cdot'}_{L^2} \in \br{\W[3/2-\s+\eps,2]}^{*} \]
   applying Proposition~\ref{prop:reg-rem} to~\eqref{eq:elg2}.
  Equation~\eqref{eq:LinearTerm} implies
  \[ \br{\tilde\rho_{k}\abs k^{-3/2+\sigma-\eps}\hat\g_{k}}_{k\in\Z} \in \ell^{2}. \]
  Recalling that $\tilde\rho_{k}\abs k^{{-3p+4}}$ converges to a positive constant as $\abs k\nearrow\infty$, we are led to
  \[ \g\in\W[\textstyle\frac{3p-4}2+\s+\frac{3p-7}2-\eps]. \]
  Choosing $\eps:=\frac{3p-7}4>0$, we gain a positive amount of regularity
  that does not depend on~$\s$. So by a simple iteration we get $\g \in W^{s,2}$
   for all $s \geq 0$.
  \end{proof}